\section{Methods}
\label{sec:methods}

\subsection{Audit design and record priority}
\label{sec:methods-audit-design}

We treated this study as an audit of existing experimental records, not as a search for a new state-of-the-art model.  Each result was linked to its protocol, configuration fingerprint, run status, prediction and target identities, metric definition, and allowed claim.  When records conflicted, final decision and summary files checked by a separately instantiated automated artifact audit took precedence over canonical result tables.  The tables took precedence over frozen protocol thresholds and historical logs.  Completed configurations were reused.  Failed or unexecuted configurations could enter later work only under a new versioned identifier, so a later run could not overwrite the earlier status.

The audit separated four questions.  \emph{Identifiability} asks whether the disclosed experiment determines a unique protocol and configuration.  \emph{Execution} asks whether that configuration reached a terminal run with an artifact.  \emph{Predictive evidence} asks whether predictions were evaluated against dataset targets under a time-respecting split.  \emph{Engineering qualification} asks whether the effect passed the minimum-effect rule and all consistency checks.  Proxy labels were used only to predict which expert would win locally.  They were not treated as forecast-accuracy evidence.

\begin{table}[t]
\centering
\caption{Compact claim-control framework. Full schemas and wording grammar are provided in the Supplementary Material.}
\label{tab:audit-framework-compact}
\footnotesize
\begin{tabular}{>{\raggedright\arraybackslash}p{0.19\linewidth}>{\raggedright\arraybackslash}p{0.36\linewidth}>{\raggedright\arraybackslash}p{0.34\linewidth}}
\toprule
Evidence object & Required control & Strongest released claim \\
\midrule
Configuration & protocol identity, disclosure mapping, execution state & registered, executed, failed, or S1--S6 evidence state \\
Prediction/target & index, time, entity, mask, and value hashes & internal/development metric for the named cell only \\
Diagnostic/proxy & target-dependent or proxy role stated explicitly & opportunity bound or predictability screen, not forecast gain \\
Engineering measure & device, precision, workload, batch, synchronization & measured latency/throughput on the frozen serving graph \\
Gate/repair & metric lock, conjunctive decision, parent/child hashes & gate pass/fail or artifact-integrity statement \\
Protected evidence & role, storage state, access definition and counter & procedural separation and zero analytical use after lock \\
\bottomrule
\end{tabular}
\end{table}


Table~\ref{tab:audit-framework-compact} summarizes the checklist used in this case.  For each claim, it records the supporting artifact, identity checks, lock or gate, and strongest allowed wording.  The checklist comes from one project and has not been validated as a general method for reducing false discoveries.  The Supplementary Material gives the full claim ledger and lock fields.

\subsection{Six-level configuration taxonomy}
\label{sec:methods-taxonomy}

We reclassified a frozen inventory of 1,224 UrbanEV configurations into six mutually exclusive evidence states.  The priority order was execution fact (S5 or S6), non-identifiability (S4), and then numerical agreement (S1--S3).

\begin{description}
    \item[S1: protocol and numerical result identifiable and reproduced.] The run covered all six folds, followed the official-compatible protocol, and each of the four disclosed metrics differed by at most 2\% from its mapped reference.
    \item[S2: protocol broadly consistent, with disclosure differences.] The run completed and remained above trend-only evidence, but used an audited correction or had a 2--5\% numerical discrepancy.
    \item[S3: relative trend only.] At least one core metric differed by more than 5\%, or the disclosed comparison cell could not be covered completely.
    \item[S4: original information insufficient for a unique configuration.] The paper omitted an identifying element, such as node identity, or an extended configuration could not be mapped uniquely to a disclosed paper cell.
    \item[S5: configuration identified but execution failed.] The manifest fixed a configuration, but the registry recorded a failed terminal state.
    \item[S6: runnable configuration not actually completed.] The manifest froze a runnable configuration that was skipped or never registered as a completed run.
\end{description}

S1--S4 can all have completed artifacts; only S1 denotes exact protocol-and-numerical reproduction under these frozen rules.  Classification was deterministic, not an inter-rater coding study.  The exact decision rules, priority order, one example per state, and reason fields are given in the Supplementary Material.  We did not change the numerical tolerances to alter the counts.  The counts describe this UrbanEV inventory and should not be read as a field-wide reproducibility rate.

\subsection{Data boundaries}
\label{sec:methods-data}

\paragraph{UrbanEV internal rolling boundary.}
We predicted hourly zonal occupancy rate; other UrbanEV quantities were outside the scope of this audit \cite{li2025urbanev}.  Occupied-port counts were divided by the frozen zone capacity.  The resulting $4{,}344\times275$ panel contained no non-finite target values and required no target imputation.  All same-protocol comparisons used a 168-hour history, horizons $\mathcal{H}=\{3,6,9,12\}$ hours, six canonical expanding-month folds, and seed 42.  Each model in a fold--horizon cell used the same target indices, timestamps, target values, and zone ordering.  The Supplementary Material lists fold ranges, counts, prediction hashes, and fit sources.  This was an internal rolling evaluation, not a blind protected test.

\paragraph{Procedurally separate Paris boundary.}
The second boundary used the Paris Smarter Mobility Data Challenge records \cite{amaraouali2024smarter,amaraouali2023parisdataset}.  For station $z$ and hour $t$, the development target was the \emph{non-available-port fraction},
\begin{equation}
    y_{t,z}=\frac{C_{t,z}+P_{t,z}+O_{t,z}}
    {A_{t,z}+C_{t,z}+P_{t,z}+O_{t,z}},
    \label{eq:paris-target}
\end{equation}
where $A,C,P,$ and $O$ denote Available, Charging, Passive, and Other counts.  Charging, Passive, and Other were all counted as unavailable because the target asks whether a port can accept a new user.  It does not measure energy demand, session volume, or active charging alone.  We locked 50 stations with stable development denominators, totaling 150 ports.  The development matrix covered 3,624 hourly timestamps from 3 July through 30 November 2020.  Of 181,200 station--hour cells, 180,073 were observed (99.378\%).

We divided the remaining timeline into a 50-day formal period (1 December 2020--19 January 2021) and a 50-day protected period (20 January--10 March 2021).  Files containing targets for these periods were stored only in a separate vault outside the main project.  We checked file identities, row boundaries, and hashes as storage metadata; the analytical-access count remained zero.  \emph{Analytical access} means using a formal or protected target for a metric, model selection, feature construction, a statistical summary, or a paper claim.  Hashing, moving, and inventorying files count only as storage operations.  Two example rows were displayed during pre-lock schema discovery, which we report as a limitation.  After the lock, no formal or protected target was used for a metric, model decision, or paper result.  The storage migration was recorded, so we do not claim that these bytes were never materialized.

Paris development models shared the same station order, targets, observation mask, history length, and horizons.  Within each fold, forward filling was limited to four hours.  Remaining missing inputs used a station-by-source-local-hour-of-week median fitted on the training interval, followed by a training-station median.  The observation mask was not a model feature, and both training loss and evaluation used observed targets only.  The released timestamps contained no UTC offsets, so we ordered the unique timestamps exactly as provided and treated adjacent entries as consecutive data hours.  The 24- and 168-hour lags therefore refer to ordinal data hours, not reconstructed UTC hours.

The exact executed model identities, local adaptations, upstream-status limits, and source hashes are provided in the Supplementary Material and in \path{models/MODEL_SOURCE_MANIFEST.csv}.

\subsection{Forecast metrics and prequential fitting}
\label{sec:methods-evaluation}

Let $\mathcal{O}_{f,h}$ be the observed target positions for fold $f$ and horizon $h$, including all valid timestamps and spatial units.  For method $m$, cell RMSE was
\begin{equation}
    R_{f,h}^{(m)}=
    \sqrt{\frac{1}{|\mathcal{O}_{f,h}|}
    \sum_{(t,z)\in\mathcal{O}_{f,h}}
    \left(\widehat y_{f,h,t,z}^{(m)}-y_{f,h,t,z}\right)^2}.
    \label{eq:cell-rmse}
\end{equation}
The primary endpoint was the equal-weight mean of cell RMSEs,
\begin{equation}
    \overline R^{(m)}=\frac{1}{|\mathcal{F}||\mathcal{H}|}
    \sum_{f\in\mathcal{F}}\sum_{h\in\mathcal{H}}R_{f,h}^{(m)}.
    \label{eq:macro-rmse}
\end{equation}
This average gives every fold--horizon cell equal weight, so cells with more observations do not dominate the main result.  MAE, WAPE, sMAPE, and RAE were retained as secondary diagnostics.  For a candidate $c$ and frozen reference $r$, positive relative gain denotes lower RMSE:
\begin{equation}
    G(c;r)=100\frac{\overline R^{(r)}-\overline R^{(c)}}
    {\overline R^{(r)}}.
    \label{eq:relative-gain}
\end{equation}
Absolute UrbanEV and Paris RMSE values were not compared because the datasets differ in spatial unit, target construction, missingness, and fold definition.

We used strict prequential fitting \cite{dawid1984prequential,tashman2000oos}.  A parameter evaluated in fold $f$ could use only validation data that preceded that fold or out-of-fold predictions from earlier folds.  For the two UrbanEV experts, CAPER ($C$) and compact TimeXer ($T$), a fixed forecast was
\begin{equation}
    \widehat y_{f,h}^{\mathrm{fixed}}
    =\alpha_{f,h}\widehat y_{f,h}^{T}
    +(1-\alpha_{f,h})\widehat y_{f,h}^{C},
    \qquad 0\leq\alpha_{f,h}\leq 1.
    \label{eq:fixed-fusion}
\end{equation}
The Global Fixed variant shared one $\alpha_f$ across horizons, whereas the Horizon Fixed control used four $\alpha_{f,h}$.  Fold 1 used only pre-test validation predictions; folds 2--6 used only earlier-fold out-of-fold predictions.  The target-dependent oracle selected the expert with smaller pointwise error.  It was an upper bound that cannot be used in practice because it observes the target.

The expert-win classifier used 21 forecast-origin features whose timestamps did not exceed the forecast origin.  Six described the experts: CAPER prediction, TimeXer prediction, signed and absolute disagreement, expert mean, and normalized horizon.  Fifteen described the observed state: current occupancy, one-step change and its magnitude, daily and weekly deviations and their magnitudes, 24-hour mean, standard deviation and turnover, normalized capacity, and sine/cosine encodings of hour and weekday.

The AUC and Brier score measured how well the classifier predicted which expert would have lower error.  They did not measure forecast accuracy or a causal mechanism.  We evaluated learned routing separately with real target RMSE.  The routing methods were a hard gate, a raw-probability soft gate, and a direct-loss multilayer perceptron.  The locked router protocol named the direct-loss MLP as the primary comparison before aggregate Gate~4 access, even though the later raw-soft point estimate was lower.

\subsection{Uncertainty estimates and progression rules}
\label{sec:methods-gates}

Paired uncertainty estimates resampled target timestamps in contiguous 24-hour blocks and kept all spatial units at a timestamp together.  The router used 1,000 bootstrap replicates.  Distillation and Paris teacher qualification used 5,000 replicates.

For Paris, we first found the timestamps shared by all four horizons in a fold and divided the ordered sequence into consecutive blocks.  Each replicate drew the original number of blocks with replacement and used the same sampled timestamps for all four horizons.  We then recomputed every cell RMSE, selected the better single-model reference inside the replicate, and calculated equal-cell macro gain.  Each monthly sequence ended with a 12-hour partial block.  Keeping this block made replicate length variable, but every original timestamp had the same expected inclusion.  We did not test 48- or 168-hour alternatives after seeing the result.  We report protocol-specific paired 95\% bootstrap intervals.  They are descriptive, and a positive interval could not override a failed minimum-effect rule.

Every listed condition had to pass.  The 1\% rule was fixed before aggregate access as a conservative project-level threshold for authorizing additional model search, cache construction, and added serving complexity.  It was not estimated from the evaluated outcomes.  It was also not derived from charger-operator utility and is not a universal significance threshold.

The lock artifacts predate their metric unlocks.  Their complete hashes and freeze/unlock times are retained in the public machine-readable audit receipt.  The displayed prefixes are router protocol \texttt{0b1fc56d...e7f18d35}, distillation protocol \texttt{fcd07f18...dfe383bc}, and Paris gate specification \texttt{96babeb1...90e9561}.

The UrbanEV router had to beat TimeXer and both fixed controls and improve macro RMSE by at least 1\%.  It also had to win at least five of six folds and 18 of 24 cells, have a paired interval lower bound above zero, and avoid more than 1\% deterioration at any horizon.  Chronos-2 teacher qualification used the same 1\%, fold, cell, and horizon criteria.  It also required matching target identity, causal input exposure, revision reproducibility, and cache feasibility.

Residual distillation compared the same 338-parameter student trained with ground truth only, aligned teacher residuals, or a 24-hour-block shuffled teacher.  Aligned KD had to improve over GT-only by at least 1\%, win five of six folds and 18 of 24 cells, have a positive paired lower bound, and keep every horizon within 1\%.  It also had to beat the shuffled control in macro RMSE, at least five folds, and the paired interval.

Paris teacher qualification was development-only.  The strict prequential CAPER--TimeXer teacher had to improve over the replicate-wise best single model by at least 1\% and win at least three of four folds and 12 of 16 fold--horizon cells.  It also needed a paired lower bound above zero, no more than 1\% deterioration at any horizon, and matching target, station, mask, timestamp, and origin identities.  A value below 1.000\% failed the point-effect condition.  We did not round the value or revise the threshold.

\subsection{Teacher, distillation, and latency settings}
\label{sec:methods-downstream}

We screened the official \texttt{amazon/chronos-2} model \cite{ansari2025chronos2} from a locked local snapshot at commit prefix \texttt{29ec3766}; the complete revision is retained in the model-lock artifact.  The screen used \texttt{chronos-forecasting 2.2.2}, zero-shot FP32 inference, a 168-hour multivariate context, no future covariates, and the native median (Q0.5) forecast.  Predictions were clipped to $[0,1]$ for the frozen primary UrbanEV RMSE, matching the existing occupancy-rate evaluation.

The residual-distillation teacher was the audited Global Fixed fusion.  The student combined a CAPER anchor with a shared DLinear residual branch \cite{zeng2023dlinear}.  GT-only, aligned-KD, and shuffled-teacher branches shared architecture, initialization, optimizer, epochs, data, and seed; only the teacher-residual term differed.  This negative control tested whether any benefit required sample-aligned teacher information rather than merely an additional regularizer.

Synchronous latency was measured on one RTX 4060 Laptop GPU in FP32 with a sequential single-device setup.  One request produced all four horizons for all 275 UrbanEV zones.  Each method--batch combination used 50 warm-up iterations and 1,000 timed iterations in each of five separately launched processes.  Batch sizes were 1, 8, and 32 origins.

Model-only time used CUDA events and synchronized the final event.  End-to-end wall-clock time synchronized the GPU before stopping.  It included input preparation, transfers, four-horizon inference, fusion, and output transfer, but excluded model loading and checkpoint I/O.  The measurement therefore represents a warm resident service, not cold start.  Batch-1 end-to-end median latency was the primary engineering endpoint.  The Supplementary Material reports the full batch 1/8/32 table and the software and hardware details.

\subsection{Fail-closed artifact audit}
\label{sec:methods-failclosed}

The pipeline kept evaluation metrics locked until every registered run closed and its predictions passed identity checks.  An invalid control mapping, target mismatch, boundary violation, incomplete provenance, or representation error blocked metric access.

Repairs were non-selective.  If one control family was defective, the full family was rerun rather than only the cells where an error was found.  Each repair receipt recorded the old and new hashes and whether predictions changed.

A separately instantiated automated artifact audit is a read-only audit created outside the workflow that generated the artifact.  It is not an external human replication, blinded coding exercise, or multi-laboratory validation.  The audit could block execution or metric access, but it could not edit predictions or lower thresholds.  Decision files were retained with negative results.  The Supplementary Material lists the states, lock fields, and unlock conditions.  The audit still relied on project-generated receipts, which remains a limitation.

The Paris storage guard also separated file storage from analytical use.  Formal and protected target-bearing files remained in vault-only storage, the main project contained no matching target source, and the analytical-access count remained zero.
